\subsection{零拷贝与 dma-heap}\label{sec:dmaheap}

一帧网球画面在送入检测网络之前要经过三个互不相同的硬件部件，先由 JPU 把 MJPEG 解码成
NV12，再由 RGA（\ref{sec:rga}）把 NV12 转换并缩放为 RGB，最后由 NPU（\ref{sec:npu}）
把 RGB 读作推理输入。这三个部件各自有独立的总线主控通道，本可以各读各的内存；倘若每一级
都把上一级的结果先搬回 CPU、再拷贝给下一级，单帧仅在两次搬运上就要付出与一次推理相当的
内存带宽与缓存抖动，感知流水线会被这些纯拷贝拖住，而它们不产生任何有用计算。让三个加速器
读写同一段物理内存、彼此之间不经过 CPU，是整条流水线能够流动起来的前提。

Linux 为这种跨设备共享提供的机制是 dma-buf 与 dma-heap。dma-heap 是一个内核内的缓冲分配器，
对应 \texttt{/dev/dma\_heap} 下的设备节点；应用向它请求一段内存，拿到的不是普通指针，而是
一个文件描述符，这个描述符代表的缓冲可以被传给任意设备去导入并直接访问。一段内存因此只分配
一次，由这个描述符在 JPU、RGA、NPU 之间传递，每个部件凭描述符解析出同一个物理地址，谁也
不必再拷贝。\textbf{本队在 Starry 上实现了这套机制，使既有的 Rockchip 用户态库无需改动即可把缓冲
在三个加速器之间流转。}

\repfig{dma-heap.png}{同一段经 \texttt{/dev/dma\_heap} 分配的物理连续缓冲以 dma-buf
描述符在 JPU、RGA、NPU 之间传递，三者就地读写而不经 CPU 拷贝}{fig:dmaheap}

\paragraph{分配器交出的是什么样的缓冲}
设备节点 \texttt{/dev/dma\_heap/system} 与 \texttt{/dev/dma\_heap/cma} 共用同一个分配器，
应用以一次 \texttt{DMA\_HEAP\_IOCTL\_ALLOC} 提交一份 \texttt{dma\_heap\_allocation\_data}
请求长度，内核返回一个可 mmap、可被各设备导入的 dma-buf 描述符。这份结构体在 LP64 下恰为
24 字节，ioctl 编号由内核侧按 \texttt{\_IOWR(\textquotesingle H\textquotesingle, 0, ...)}
推算并以编译期断言固定为 \texttt{0xC0184800}，与主线 Linux 的取值逐位一致，任何字段或编号的
偏移都会在编译时暴露，使用户态库据以组包的布局不会悄然漂移。分配本身落在 4\,GiB 以下的
物理空间，按页（4\,KiB）对齐且物理连续。这两条约束并非随意而来。物理连续是因为 RGA2 与 NPU
在本阶段以设备侧地址翻译部件（IOMMU）旁路、物理直通的方式工作，它们沿一个起始物理地址顺序
访问整段缓冲，中途若被分页打断便会读到错误的页；32 位地址上限则对应 RGA2 把一个裸 32 位
物理基址直接写入自身寄存器的约束，因而分配器以 \texttt{u32::MAX} 作为 DMA 掩码，并把请求
长度限制在 32 位以内。

\paragraph{同一段内存上的三个读者}
JPU 把解码出的 NV12 写入这段缓冲，RGA 以 NV12 读入、RGB 写出，NPU 再把 RGB 读作输入，
全程没有一次 CPU 拷贝。串起这三者的是同一个 dma-buf 描述符，每个部件通过导入它而拿到那个
唯一的物理地址。RGA 一侧的导入入口在 \texttt{/dev/rga} 收到一个描述符或一个先前登记的句柄，
经 \texttt{resolve\_dmabuf\_fd} 把它解析回那段分配，取出物理地址交给引擎。关键在于这次解析
同时克隆了该分配的引用计数，并把这份克隆在整个提交与轮询期间持有；缓冲的生命周期由这个引用
计数管理，最后一个持有者释放时页面才归还。这样即便另一线程在引擎运行途中并发地释放同一缓冲，
内存也不会被中途抽走，引擎读写的始终是有效的页。

\paragraph{让 CPU 与设备看到同一份数据}
设备物理直通带来一个缓存一致性问题。这段连续缓冲在 aarch64 上是带缓存的，分配时由 CPU 清零，
那些清零的缓存行处于脏状态；若不先把它们写回内存，引擎读到的将是旧值，或一次稍后的缓存逐出
会覆盖掉引擎刚写出的结果。因此在把缓冲交给引擎之前，沿写入设备的方向显式同步一次，将脏的
CPU 缓存行刷回内存；待引擎完成、在读取结果之前，再沿读回 CPU 的方向同步一次，使目标缓冲对应
的缓存行失效，CPU 随后读到的是引擎真正写入内存的数据而非陈旧的缓存。这两次方向相反的同步
分别在引擎起始之前与轮询到完成之后进行，只针对实际参与本次操作的源与目标缓冲，CPU 与设备
据此看到一致的内容。

这一层是 RGA（\ref{sec:rga}）、JPU（\ref{sec:jpu}）与 NPU（\ref{sec:npu}）三节共同依赖的
底座，它把原本散落在各级之间的拷贝从流水线上彻底移除，使一帧只被分配一次、被三个加速器依次
就地读写。
